\chapter{総合演習：小さな作品を完成させる}
\label{chap:final-projects}

この章では、これまでに学んだ内容を使って、1つの作品を完成させます。ここでの目標は、すべての技術を使うことではありません。自分で課題を選び、完成までの道筋を考え、必要な機能を少しずつ作っていく経験をすることです。

総合演習では、学生ごとに取り組む課題が違っていてかまいません。得意な人は発展的な機能に挑戦し、まだ自信がない人は、まず確実に動く小さな作品を作ることを目指します。作品の大きさよりも、「自分で考えて、動く形にした」ことを大切にします。

\section{この章の概要}

この章は、通常の解説章ではなく、制作課題を選ぶためのガイドです。各課題は、1週間から2週間程度で取り組むことを想定しています。

この章で扱う内容は次の通りです。

\begin{itemize}
    \item 総合演習の進め方
    \item 課題を選ぶときの考え方
    \item 提出物の条件
    \item 評価の観点
    \item 難易度別の課題一覧
    \item 制作計画と振り返り
\end{itemize}

\begin{pointbox}{完成の大きさを自分で調整する}
総合演習では、同じ課題でも完成の形を変えられます。まず最低限の完成条件を満たし、時間があれば見た目、音、データ保存、難易度調整などを追加していきましょう。
\end{pointbox}

\section{総合演習の進め方}

総合演習では、課題一覧から1つを選びます。選んだ課題をそのまま作ってもよいですし、テーマや見た目を自分なりに変えてもかまいません。

制作は、次のような流れで進めます。

\begin{enumerate}
    \item 課題を1つ選ぶ
    \item 最低限の完成条件を確認する
    \item 画面に何を表示するかを紙やメモに描く
    \item 必要な変数、関数、クラスを考える
    \item まず小さく動くものを作る
    \item 少しずつ機能を追加する
    \item 動作確認をして、説明できる形に整える
\end{enumerate}

\begin{notebox}{最初から完成形を作ろうとしない}
最初からすべての機能を入れようとすると、どこで動かなくなったのか分かりにくくなります。まずは「丸が動く」「点数が増える」「画面が切り替わる」のような小さな成功を作りましょう。
\end{notebox}

\section{提出物の条件}

提出物は、作品そのものだけではなく、作品を説明するための情報も含めます。プログラムを見た人が、何を作ろうとしたのか、どこを工夫したのかを理解できるようにします。

提出物には、次の内容を含めます。

\begin{itemize}
    \item TypeScript + p5.jsで作成したプログラム
    \item 作品名
    \item 作品の説明
    \item 操作方法
    \item 工夫した点
    \item うまくいかなかった点や、今後追加したい点
\end{itemize}

プログラムは、授業で扱ってきたp5.jsのinstance modeを基本とします。変数名、関数名、クラス名は、あとから読んでも意味が分かる名前にします。

\section{評価の観点}

総合演習では、単に機能が多い作品だけを高く評価するのではありません。自分の理解に合わせて設計し、読みやすく動く形にしたことも重要です。

評価では、主に次の観点を見ます。

\begin{center}
\begin{tabular}{ll}
\toprule
観点 & 内容 \\
\midrule
動作 & 作品が実行でき、主要な機能が動く \\
読みやすさ & 変数名、関数名、クラス名が分かりやすい \\
構造 & 関数、配列、クラスなどを適切に使っている \\
工夫 & 見た目、動き、ルール、演出に自分なりの工夫がある \\
説明 & 作品の目的や操作方法を説明できる \\
振り返り & できたこと、難しかったこと、改善点を書ける \\
\bottomrule
\end{tabular}
\end{center}

\begin{pointbox}{難しい機能を使えばよいわけではない}
高度な機能を無理に入れるより、作品の目的に合った書き方を選ぶことが大切です。小さくても、よく整理されていて遊びやすい作品はよい作品です。
\end{pointbox}

\section{難易度の目安}

課題は、AからDまでの難易度に分けています。Aは基礎を確認しながら作る課題、Dは発展的な設計に挑戦する課題です。

\begin{center}
\begin{tabular}{lll}
\toprule
難易度 & 目安 & 主に使う内容 \\
\midrule
A & 基礎 & 図形、変数、条件分岐、マウス、キーボード \\
B & 標準 & 配列、関数、クラス、当たり判定、ゲームの状態 \\
C & 発展 & 画像、文字、音声、JSON、保存、アルゴリズム \\
D & 挑戦 & 継承、interface、設計の工夫、自由制作 \\
\bottomrule
\end{tabular}
\end{center}

自分にとって少し難しいと感じる課題を選ぶと、制作を通して学べることが増えます。ただし、難しすぎる課題を選ぶと、最後まで動く形にできないことがあります。迷ったときは、最低限の完成条件を見て選びましょう。

\section{A：基礎を固める課題}

このレベルの課題では、まず1つの画面で動く作品を作ることを目指します。図形、色、変数、条件分岐、マウスやキーボード操作を使って、確実に動く作品にします。

\subsection{A1 キーボードで動かすキャラクター}

画面上のキャラクターをキーボードで動かす作品を作ります。キャラクターは円や四角形だけで作ってもかまいません。

\begin{itemize}
    \item 使う主な内容：変数、条件分岐、キーボード入力、図形描画
    \item 最低限の完成条件：矢印キーまたはWASDキーでキャラクターを動かせる
    \item 追加してよい工夫：画面端で止まる、背景を描く、表情を変える、ゴールを作る
    \item 評価の観点：操作が分かりやすく、位置を表す変数名が読みやすい
\end{itemize}

\subsection{A2 マウスで遊ぶお絵描きツール}

マウスを使って画面に線や図形を描くツールを作ります。色や太さを変えられるようにすると、作品としての楽しさが増します。

\begin{itemize}
    \item 使う主な内容：マウス入力、色、条件分岐、変数
    \item 最低限の完成条件：マウスをドラッグすると線や図形を描ける
    \item 追加してよい工夫：キーで色を変える、太さを変える、画面を消す機能を作る
    \item 評価の観点：操作方法が明確で、描画処理が読みやすく整理されている
\end{itemize}

\subsection{A3 図形で作る時計やタイマー}

円、線、文字を使って、時計やタイマーのような表示を作ります。正確な時計でなくても、時間の変化が見て分かる作品であればよいです。

\begin{itemize}
    \item 使う主な内容：変数、文字表示、角度、アニメーション
    \item 最低限の完成条件：時間の経過に合わせて表示が変化する
    \item 追加してよい工夫：針を回す、残り時間を表示する、終了時に色を変える
    \item 評価の観点：数値と見た目の対応が説明できる
\end{itemize}

\subsection{A4 クリックで反応するインタラクティブ作品}

クリックした場所に図形を出したり、クリックするたびに画面の様子が変わったりする作品を作ります。ゲームでなくても、触って変化する作品であればかまいません。

\begin{itemize}
    \item 使う主な内容：マウスクリック、条件分岐、配列、色
    \item 最低限の完成条件：クリックに反応して画面が変化する
    \item 追加してよい工夫：クリックした図形が動く、一定時間で消える、音や文字を加える
    \item 評価の観点：反応が分かりやすく、状態を表す変数が整理されている
\end{itemize}

\section{B：標準的な制作課題}

このレベルの課題では、複数のものを動かしたり、ゲームのルールを作ったりします。配列、関数、クラスを使うと、作品を整理しやすくなります。

\subsection{B1 アイテムを集めるミニゲーム}

プレイヤーを動かし、画面上に出てくるアイテムを集めるゲームを作ります。点数や制限時間を入れると、ゲームとして分かりやすくなります。

\begin{itemize}
    \item 使う主な内容：キーボード入力、配列、当たり判定、点数
    \item 最低限の完成条件：プレイヤーがアイテムに触れると点数が増える
    \item 追加してよい工夫：制限時間、減点アイテム、難易度の上昇、効果音
    \item 評価の観点：当たり判定と点数処理が関数などで読みやすく分けられている
\end{itemize}

\subsection{B2 障害物をよけるゲーム}

上から落ちてくる障害物や、横から流れてくる障害物をよけるゲームを作ります。プレイヤーが障害物に当たるとゲーム終了になるようにします。

\begin{itemize}
    \item 使う主な内容：アニメーション、配列、条件分岐、ゲームの状態
    \item 最低限の完成条件：障害物が動き、当たるとゲームオーバーになる
    \item 追加してよい工夫：スタート画面、リトライ、スコア、画像の利用
    \item 評価の観点：ゲーム中、ゲームオーバーなどの状態が分かりやすく管理されている
\end{itemize}

\subsection{B3 複数のボールや敵が動く作品}

複数のボール、敵、星、粒子などが画面上を動く作品を作ります。動きの違いを作ると、画面に変化が出ます。

\begin{itemize}
    \item 使う主な内容：配列、繰り返し、関数、クラス
    \item 最低限の完成条件：複数のものが同時に動く
    \item 追加してよい工夫：跳ね返り、重力のような動き、色や大きさの違い、クリックで追加
    \item 評価の観点：1つのものを表すデータやクラスが整理されている
\end{itemize}

\subsection{B4 画面が切り替わるクイズアプリ}

問題文と選択肢を表示し、回答に応じて正解や不正解を表示するクイズアプリを作ります。ゲーム制作が苦手な人でも取り組みやすい課題です。

\begin{itemize}
    \item 使う主な内容：文字表示、配列、条件分岐、画面状態
    \item 最低限の完成条件：複数の問題を順番に表示し、正解数を数えられる
    \item 追加してよい工夫：画像つき問題、制限時間、結果画面、問題のランダム表示
    \item 評価の観点：問題データと表示処理が分かれている
\end{itemize}

\section{C：発展的な制作課題}

このレベルの課題では、画像、音声、ファイル入出力、アルゴリズムの工夫などを使います。基本機能を先に完成させてから、発展機能を追加していくと進めやすくなります。

\subsection{C1 JSONで設定を読み込むゲーム}

ゲームの初期設定をJSONから読み込む作品を作ります。敵の数、アイテムの種類、制限時間などをデータとして外に出すことで、プログラムを変えずに調整できるようにします。

\begin{itemize}
    \item 使う主な内容：JSON、オブジェクト、配列、ファイル読み込み
    \item 最低限の完成条件：JSONに書いた設定を使ってゲームを開始できる
    \item 追加してよい工夫：複数ステージ、難易度選択、アイテム設定の追加
    \item 評価の観点：JSONの構造とTypeScript側のデータの対応を説明できる
\end{itemize}

\subsection{C2 ハイスコアを保存するゲーム}

ゲームの得点を保存し、次に実行したときにもハイスコアを表示できる作品を作ります。スコアの保存は、ゲームらしさを高めるよい題材です。

\begin{itemize}
    \item 使う主な内容：得点、ファイル入出力、文字表示、ゲーム状態
    \item 最低限の完成条件：ゲーム終了時に得点を保存し、ハイスコアを表示できる
    \item 追加してよい工夫：プレイヤー名、ランキング、日付、難易度別スコア
    \item 評価の観点：保存するデータが必要最小限に整理されている
\end{itemize}

\subsection{C3 音声効果つきミニゲーム}

アイテムを取ったとき、失敗したとき、ゲームが始まったときなどに音を鳴らすミニゲームを作ります。音は多すぎると分かりにくくなるため、場面に合った使い方を考えます。

\begin{itemize}
    \item 使う主な内容：音声、条件分岐、ゲーム状態、画像や図形
    \item 最低限の完成条件：少なくとも2種類の場面で音が鳴る
    \item 追加してよい工夫：BGM、音量調整、効果音の使い分け、タイトル画面
    \item 評価の観点：音がゲームの状態や操作結果と対応している
\end{itemize}

\subsection{C4 ライフゲームの改造}

第\ref{chap:algorithm-improvement}章で扱ったライフゲームをもとに、ルールや見た目、操作方法を変えた作品を作ります。速く動かす工夫に挑戦してもよいです。

\begin{itemize}
    \item 使う主な内容：2次元配列、繰り返し、条件分岐、アルゴリズム
    \item 最低限の完成条件：セルの状態が世代ごとに変化する
    \item 追加してよい工夫：初期配置の編集、色の変化、速度比較、独自ルール
    \item 評価の観点：状態の更新方法を説明でき、前の世代と次の世代を混同していない
\end{itemize}

\section{D：挑戦的な制作課題}

このレベルの課題では、プログラムの設計そのものを考える力が必要になります。完成した機能の数だけでなく、どのように整理したか、どこまで試したかも大切です。

\subsection{D1 RPG風バトル}

勇者、魔法使い、モンスターなどのキャラクターをクラスで表し、簡単なバトルを作ります。継承を使って共通部分をまとめると、第14章の内容を発展させられます。

\begin{itemize}
    \item 使う主な内容：クラス、継承、メソッド、HP、攻撃処理
    \item 最低限の完成条件：2種類以上のキャラクターが攻撃し、HPが変化する
    \item 追加してよい工夫：魔法、回復、敵の種類、ターン制、画像や音
    \item 評価の観点：共通する性質と違う動きがクラスで整理されている
\end{itemize}

\subsection{D2 interfaceを使った収集ゲームの拡張}

第\ref{chap:polymorphism-leaf-collector}章の考え方を使い、複数種類のアイテムを同じ配列で扱うゲームを作ります。見た目や効果が違うものを、同じ呼び方で扱えるようにします。

\begin{itemize}
    \item 使う主な内容：interface、クラス、配列、ポリモーフィズム
    \item 最低限の完成条件：2種類以上のアイテムを同じ配列で更新・描画できる
    \item 追加してよい工夫：回復アイテム、減点アイテム、時間延長、レアアイテム
    \item 評価の観点：interfaceを使う理由を、自分の作品に合わせて説明できる
\end{itemize}

\subsection{D3 自作ルールのシミュレーション}

人、粒子、生き物、交通、天気など、自分で決めたルールにしたがって変化するシミュレーションを作ります。正確な科学シミュレーションである必要はありません。

\begin{itemize}
    \item 使う主な内容：配列、クラス、条件分岐、アルゴリズム
    \item 最低限の完成条件：複数の要素がルールにしたがって変化する
    \item 追加してよい工夫：パラメータ変更、グラフ表示、速度比較、ランダム性
    \item 評価の観点：ルールとプログラムの対応を説明できる
\end{itemize}

\subsection{D4 自由制作}

ここまでの課題に当てはまらない作品を作りたい場合は、自由制作として取り組めます。ただし、自由制作でも最低限の完成条件と評価の観点を自分で決めてから始めます。

\begin{itemize}
    \item 使う主な内容：作品に応じて選ぶ
    \item 最低限の完成条件：自分で決めた目的を満たす動く作品にする
    \item 追加してよい工夫：授業で扱った内容を組み合わせる、見た目や操作性を整える
    \item 評価の観点：何を作り、どこを工夫したのかを説明できる
\end{itemize}

\begin{notebox}{自由制作は計画を小さくする}
自由制作では、やりたいことが大きくなりがちです。最初に「最低限ここまでできれば完成」と言える形を決めてから始めましょう。
\end{notebox}

\section{制作計画を立てる}

課題を選んだら、すぐにコードを書き始める前に、制作計画を立てます。計画は長く書く必要はありませんが、何から作るかを決めておくと、途中で迷いにくくなります。

制作計画には、次の内容を書きます。

\begin{itemize}
    \item 選んだ課題名
    \item 作品名
    \item 画面に表示するもの
    \item 操作方法
    \item 最低限完成させる機能
    \item 余裕があれば追加したい機能
    \item 使う予定の変数、関数、クラス
\end{itemize}

\begin{pointbox}{最初の目標は小さくする}
「完成したら楽しい作品」を考えることは大切ですが、最初の目標は小さくします。たとえば、収集ゲームなら、最初はプレイヤーと1個のアイテムだけで当たり判定を作ります。
\end{pointbox}

\section{制作中に確認すること}

制作中は、プログラムが大きくなりすぎる前に、少しずつ確認します。動かなくなってから一気に直すより、動いている状態を保ちながら進める方が安心です。

次の点を、ときどき確認しましょう。

\begin{itemize}
    \item 変数名は意味が分かるか
    \item 同じ処理を何度も書いていないか
    \item 関数に分けた方がよい処理はないか
    \item クラスにした方がよいものはないか
    \item コメントは多すぎず、少なすぎないか
    \item ゲームの状態が分かりやすく管理されているか
    \item 作品の操作方法を他の人に説明できるか
\end{itemize}

\section{発表と振り返り}

作品が完成したら、動かして見せるだけでなく、どのように作ったかを説明します。うまくいったところだけでなく、難しかったところを話すことも大切です。

発表では、次の内容を説明できるようにします。

\begin{itemize}
    \item どの課題を選んだか
    \item 作品の目的や遊び方
    \item 工夫したところ
    \item プログラムで特に見てほしいところ
    \item 難しかったところ
    \item さらに時間があれば追加したいところ
\end{itemize}

振り返りでは、完成度だけでなく、自分がどのように考えて作ったかを書きます。制作の途中で予定を変えた場合も、その理由を書ければ十分に意味があります。

\section{まとめ}

この章では、総合演習として1つの作品を完成させるための課題を紹介しました。課題は難易度別に分かれていますが、どの課題を選んでも、まずは小さく動く形を作ることが大切です。

プログラミングでは、最初から大きな作品を一度に作ることはほとんどありません。小さく作り、動かし、直し、少しずつ広げることで、作品は完成に近づいていきます。この教材で学んだ内容を組み合わせ、自分なりの作品として形にしてみましょう。

\section{演習問題}

この章の演習問題は、総合演習を進めるための準備と振り返りです。選んだ制作課題と合わせて取り組みましょう。

\begin{enumerate}
    \item 課題一覧から、自分が取り組みたい課題を1つ選び、その理由を書きなさい。
    \item 選んだ課題の最低限の完成条件を、自分の言葉で書き直しなさい。
    \item 作品の画面案を1つ描き、どこに何を表示するか説明しなさい。
    \item 作品で使う予定の変数を5個以上挙げ、それぞれ何を表すか説明しなさい。
    \item 作品で使う予定の関数を3個以上挙げ、それぞれの役割を書きなさい。
    \item クラスを使う場合は、どのようなものをクラスとして表すか説明しなさい。
    \item 作品に追加したい工夫を3つ挙げ、優先順位を付けなさい。
    \item 途中で動かなくなったときに、どのように原因を調べるかを書きなさい。
    \item 完成後、作品の操作方法を初めて使う人向けに説明しなさい。
    \item 制作を終えて、できるようになったこと、難しかったこと、次に挑戦したいことを書きなさい。
\end{enumerate}
